\chapter{Nephritis}

\medskip
\noindent\textit{\textbf{Under review.} This chapter is an AI-assisted educational draft; it is not a substitute for professional medical advice.}

\section*{Definition and Overview}
inflammation of the kidneys. It may affect the glomeruli, tubules, or surrounding tissue and can impair kidney function. Individual assessment is important because symptoms and treatment vary with the cause.

\section*{Clinical Features}
Blood or protein in urine, swelling, high blood pressure, reduced urine, flank pain, fever, or fatigue.

\section*{When to Seek Medical Care}
Seek medical review for new, persistent, worsening, or function-limiting symptoms. Seek emergency help for severe breathing difficulty, collapse, sudden neurological change, severe chest or abdominal pain, or any rapidly worsening illness.

\section*{Diagnosis and Investigations}
Urinalysis and urine protein measurement, kidney-function blood tests, blood pressure, immune and infection tests, ultrasound, and sometimes kidney biopsy.

\section*{Management}
Treat the cause, control blood pressure and fluid balance, stop harmful medicines when advised, and use specialist immunosuppression only when indicated.

\section*{Complications}
Possible complications depend on the cause and may include progression of the condition, effects on other organs, treatment adverse effects, and reduced quality of life. Early assessment can reduce preventable harm.

\section*{Prognosis and Outcomes}
Outcome depends on the underlying cause, severity, complications, and response to treatment. Discuss individual prognosis with the treating clinician.

\section*{Prevention and Self-Care}
Follow the treatment plan, keep relevant appointments, avoid known triggers, protect affected areas from injury, and seek help early if symptoms worsen. Do not start or stop prescription treatment without clinical advice.

\section*{Expanded clinical context}
Nephritis is a infection-related topic. Interpretation should account for age, baseline health, medicines, comorbidities, goals, and access to care. This chapter supports discussion with a qualified clinician and does not establish a diagnosis.

\section*{Assessment and decision points}
Assessment may include examination, exposure and travel history, immune status, targeted swabs or cultures, molecular tests, blood tests, or imaging. Not every infection requires antibiotics. Document the main concern, time course, functional impact, relevant risks, and red flags. Use targeted investigations when their result is likely to change management.

\section*{Management and follow-up}
Treatment depends on the organism and site and may include fluids, rest, isolation advice, wound care, antimicrobials, and management of complications. Explain expected improvement, when reassessment is needed, and how follow-up can be accessed. Avoid starting, stopping, or sharing prescription medicines without professional advice.

\section*{Safety-netting}
Seek urgent care for sepsis features, breathing difficulty, confusion, severe dehydration, meningitis symptoms, or rapid deterioration in a vulnerable person.

\section*{Practical prevention}
Use vaccination, hand hygiene, food and water safety, protective equipment, and current public-health advice. Keep written instructions and seek review if the topic recurs, changes, or does not improve as expected.

\section*{Limitations}
This educational expansion should be checked against current local guidance and authoritative topic-specific sources.
\clearpage
